\section{Границы метода и открытые вопросы}\label{sec:open}

\begin{enumerate}[leftmargin=1.5em]
\item \textbf{Насколько мал минимальный индекс?} Порог $d=1$ пробит вплоть
до индекса~$43$ (\S\ref{ssec:constr43}). История этого спуска поучительна.
Кампания 2026-07 оптимизировала форму Грама по всем девяти существенным
параметрам (CMA-ES, восемь независимых инстансов $\times\,2200$ точных оценок
на $k=44$, мультистарты Нелдера--Мида от найденных оптимумов) и вышла на плато
$d=0{,}990$ при $k=44$ и $0{,}971$ при $k=43$; вдоль лестницы
$k=43,\dots,31$ наилучшая найденная ширина падала до $0{,}849$. Дополнительно
для проверенных пар <<решётка\,--\,структура факторгруппы>> при $36\le k\le44$
минимум числа неотделённых запрещённых векторов равнялся~$1$. Всё это было
честно записано как отрицательный экран~--- и оказалось \emph{недобором
поиска}: у $k=43$ было всего три рестарта, и ни один не попал в область, где
живёт эйзенштейнов максимум $1{,}00411$. Мораль методическая: сужение
пространства симметрией (\S\ref{ssec:symred}) оказалось сильнее наращивания
рестартов в полном пространстве.

Ниже $43$ порог не пройден. Проверено (статус~\textbf{[Э]}, артефакты ---
\path{audit-data/results/dim4_below43_screen.json},
\path{audit-data/results/dim4_eisenstein_scan.json},
\path{audit-data/results/dim4_cyclotomic_scan.json},
\path{audit-data/results/dim4_symmetry_atlas.json}):
\begin{itemize}[leftmargin=1.3em,topsep=2pt]
\item эйзенштейново трёхпараметрическое семейство --- плотная сетка по
приведённой области ($\det h=1$, $a\le b$, $|c|\le a/\sqrt3$; касповый конец
$a<0{,}4$, где решётка вытягивается, а $d$ падает, отсечён; $606$ решёток), и
для \emph{каждой} решётки --- штатный исчерпывающий перебор всех эрмитовых
нормальных форм индекса $k$ с локальной доводкой. Итог при $31\le k\le42$:
$0{,}776\le d\le0{,}951$, максимум $0{,}9507$ при $k=39$; при $k=43$ тот же
перебор даёт $1{,}00411$. Неинвариантные $\Gamma$ здесь покрыты наравне с
инвариантными;
\item циклотомические семейства ранга~1 ($\Z[\zeta_5]$, $\Z[\zeta_8]$,
$\Z[\zeta_{12}]$; один параметр) --- максимум $d=0{,}9275$, и тот при $k=44$;
\item инвариантные подмодули эйзенштейнова семейства для \emph{всех}
норменных индексов $k\le43$ --- лестница
\[
\begin{array}{c|ccccccccccccc}
k & 7&9&12&13&16&19&21&25&27&28&31&36&37\\\hline
d & 0{,}281&0{,}424&0{,}480&0{,}523&0{,}633&0{,}628&0{,}670&0{,}714&
0{,}698&0{,}827&0{,}840&0{,}889&0{,}934
\end{array}
\qquad
\begin{array}{c|cc}
k & 39 & \mathbf{43}\\\hline
d & 0{,}951 & \mathbf{1{,}004}
\end{array}
\]
то есть внутри эйзенштейнова класса порог $d=1$ впервые достигается ровно при
$k=43$ --- первой норме, на которой это возможно;
\item атлас классов $\Z[S]$-модулей ранга~4 (все с инвариантным
пространством не выше четырёх параметров) --- случайный поиск по инвариантному
конусу с доводкой и тем же исчерпывающим перебором подрешёток: гауссово
семейство ($\Z[i]$-модули ранга~2) даёт при $k\le42$ не более $0{,}865$,
модули $\Z[\omega]\oplus\Z^2$ --- не более $0{,}802$, а \emph{склеенный}
модуль $\Z[C_3]\oplus\Z$ (регулярное представление $C_3$ плюс тривиальное
слагаемое; в прямую сумму подрешёток не разлагается) --- $0{,}9583$ при
$k=42$. Последнее число само по себе рекорд для этого индекса: прежний
чемпион кампании 2026-07, найденный поиском по всем формам, давал $0{,}9489$.
Порог не взят и здесь, но видно, что симметрийное сужение улучшает не только
пробитый индекс~$43$, но и ближайшие непробитые.
\end{itemize}
Замечание о том, куда целиться дальше. Индексы $40$, $41$, $42$ выпали из
инвариантного перебора не по геометрической, а по \emph{арифметической}
причине: $42$ не является нормой ни в $\Z[\omega]$, ни в $\Z[i]$, поэтому
подрешётки нужной формы там просто не существует (то же квантование, что в
\S\ref{ssec:quant} для плоскости). А мозаичная схема
определения~\ref{def:tiling} этого ограничения не знает: в плоскости именно на
нелёшиановых индексах она обгоняет решёточную. Пол мозаичного принципа
(теорема~\ref{thm:floor}) даёт в $\R^4$ лишь $2^4=16$, так что $k=42$ им не
запрещён. Это и есть самая мотивированная следующая попытка --- мозаика на
$40$--$42$ цвета в $\R^4$.

Как и прежде, всё сказанное --- \emph{не} доказательство минимальности:
нижняя граница Кулсона $\chi_s\ge2^{n+1}-1=31$ \cite{Coulson,ABPR}
формального препятствия не даёт, а только что описанная история $45\to43$ --- второй после индекса
$134$ в $\R^5$ случай, когда <<экран>> снимался переходом в другую область, а
не уточнением прежней. Строгое доказательство минимальности (или конструкция
с $31\le k\le42$) --- открытая задача в духе вопроса~1 из \cite{ABPR}.
\item \textbf{Аналитическое описание.} Найти замкнутую форму оптимальных
решёток для индексов $43$, $45$ и $48$ (предельных точек максимизации $d$).
Для $43$ задача выглядит доступнее прочих: оптимум лежит в
трёхпараметрическом эйзенштейновом семействе и характеризуется равенством
трёх орбит векторов $\Gamma$ (замечание после теоремы~\ref{thm:main43}), то
есть является решением системы из трёх алгебраических уравнений на три
неизвестных.
\item \textbf{Размерность 5: ниже $132$.} Предел Кулсона здесь равен $63$, то
есть формального препятствия к дальнейшему спуску нет. Тот же трёхшаговый
механизм (полный поиск ядра, деформация формы, доводка по активному набору),
применённый к индексам $131$, $130$ и $129$, порога не пересёк: лучшие
отношения равны соответственно $0{,}990216$, $0{,}976688$ и
$0{,}957014$~\textbf{[Э]}. Для каждого из них исчерпывающий конечнополевой
оракул подтвердил, что это глобальный дискретный максимум \emph{на
соответствующей фиксированной численной форме}; подробности ---
приложение~\ref{app:d5}.

\item \textbf{Размерность 6: оценка $343$ не улучшена.} Предел Кулсона
равен~$127$. Мы проверили как ближайшие индексы, так и арифметически
согласованные цели; ни один не достиг порога~\textbf{[Э]}:

\smallskip
\begin{center}\small
\begin{tabular}{c l l}
\toprule
$k$ & структура факторгруппы & лучшее $D/\diam V_0$\\
\midrule
$342=19\cdot3\cdot3\cdot2$ & циклическая / нециклическая & $0{,}940393$ / $0{,}924370$\\
$341$ & --- & ${\approx}\,0{,}9164$\\
$340=17\cdot5\cdot2\cdot2$ & циклическая / нециклическая & $0{,}935190$ / $0{,}928955$\\
$338=13\cdot13\cdot2$ & нециклическая & $0{,}908895$\\
$333=37\cdot3\cdot3$ & нециклическая & $0{,}961088$\\
$329$, $322$, $315$ & сохраняют исх. характер & $0{,}944303$, $0{,}949097$, $0{,}933317$\\
$\mathbf{336=7\cdot4\cdot4\cdot3}$ & Smith $(1,1,1,1,4,84)$ & $\mathbf{0{,}984244}$\\
\bottomrule
\end{tabular}
\end{center}
\smallskip

Наиболее перспективна не ближайшая, а арифметически согласованная цель $336$:
фактор точной конструкции $343$ имеет Smith-тип $(1,1,1,7,7,7)$, и, сохранив
один исходный характер над $\mathbb F_7$, мы ищем три остаточных блока
совместно с формой. Точный переход через стены L-типа, последовательный SDP с отсекающими плоскостями и внешняя аппроксимация конуса линейными сечениями довели
отношение до $0{,}984244122$, то есть разрыв до порога сократился с $1{,}934\%$
до $1{,}576\%$. Три независимых подхода упираются в один и тот же барьер, что
указывает на исчерпанность текущего вторичного конуса: дальше нужно менять сам
конус, согласуя выбор периода с выбором формы. Проверены также раскраски, не
являющиеся раскрасками по классам смежности (периодические подъёмы, графы Кэли
конечного тора, аффинные семьи); ни одна не дала раскраски. Все детали ---
приложение~\ref{app:d6}.

\item \textbf{Размерность 7.} Примарный Radon-поиск и деформация формы понизили
оценку АБПР с $1372$ до $1323$ (теорема~\ref{thm:r7}), хотя все прямые атаки на
фиксированную $E_7^*$ давали барьер. Индексы $1322$, $1321$ и $1320$ проверены
как шаги $-1,-2,-3$; лучший дискретно-непрерывный фронтир получен при $1320$
и равен $0{,}975568590$. Повторный локальный поиск ядра на этой форме
воспроизвёл исходное ядро, но порог не пересёк. Следующие задачи --- продолжить
из нескольких ядер и добавить совместный штраф доверительной области. Разрыв с нижней
границей всё ещё огромен: $16\le\chi(\R^7)\le1323$.
\item \textbf{Размерности 8 и 9.} Ближайшие экраны при $2400$ (для $E_8$) и
$17\,250$ (для $A_9^*$) дают лишь $0{,}76425$ (после деформации формы;
$0{,}70710678$ до неё) и $0{,}825837653$~\textbf{[Э]};
оценки $2401$ и $17\,253$ не изменяются. Smith-структура
$\operatorname{diag}(1^4,7^4)$ объясняет, почему малые поправки здесь особенно
неэффективны: возмущение ранга меньше четырёх меняет определитель на число,
кратное~$7$, и потому не может дать $2400=2401-1$. Детали и найденная в
\cite{ABPR} смена координат для $C_9$ --- приложение~\ref{app:d89}.

\item \textbf{Тождество~\eqref{eq:eis} --- закрыто.} В предыдущей версии здесь
стоял открытый вопрос: доказано было лишь неравенство
$D((3+\omega)\Lambda)^2\ge\tfrac73\lambda_1^2$, а равенство проверено численно.
Теперь равенство доказано для всех эйзенштейновых решёток
(теорема~\ref{thm:eis}); шестиугольное сечение ячейки действительно никогда не
срезается (следствие~\ref{cor:hexcell}). Открытым остаётся смежный вопрос:
описать все решётки с <<$A_2$-треугольником>> минимальных векторов, для которых
плоское сечение $V_0\cap P=H$ даёт \emph{глобальный} минимум $D$, а не только
минимум по векторам плоскости $P$ (замечание~\ref{rem:a2tri}) --- то есть когда
экран $d\le\sqrt{7/3}/\rho$ достигается.

\item \textbf{Куда целиться: открытые окна в числах.} Соберём в одном месте
цели, на которые указывают экраны и учёт этой работы.
\begin{itemize}[leftmargin=1.2em,itemsep=1pt]
\item \emph{$K_{12}$ --- самый большой запас} (\S\ref{ssec:window}):
подрешётка $\Gamma\subset K_{12}$, улучшающая $3^{12}$, обязана иметь
$\lambda_1^2\ge28$ и индекс в окне $[111\,019,\ 3^{12}]$ --- свободный
множитель $4{,}79$; лестница подобий в окне пуста, целиться надо в
$\Z[\omega]$-подмодули с произвольной нормой определителя.
\item \emph{Разрыв восьмимерной лестницы} $2401$ ($d=1{,}0801$) --- $6561$
($d=1{,}4142$) --- общее узкое место размерностей $10$--$16$
(\S\ref{ssec:prodneg}): раскраска $\R^8$ индекса ${\approx}2800$--$3800$ с
$d\approx1{,}15$--$1{,}25$ немедленно дала бы слой ранга~$2$ и
${\approx}243\,000$ в $\R^{12}$. Деформация $E_8$, по-видимому, закрыта
(локальная минимальность $\rho$, \S\ref{ssec:prodneg}); остаются
неэйзенштейновы подрешётки.
\item \emph{Плоские индексы $17$, $18$} исключены лишь численно --- и для
решёток (\S\ref{ssec:planefloor}), и для мозаик (\S\ref{ssec:quant});
доказательный пол стоит на $16{,}021$.
\item \emph{Рационализация} оставшихся граничных сертификатов: $21609$ в
$\R^{10}$ (запас $4{,}8\cdot10^{-4}$, \S\ref{ssec:dim10lam}) и перевод
ламинирований $7203$/$28812$ из \textbf{[Ч]} в \textbf{[С]}; это вынесет их в
заголовок вместо нынешних $9604$ и $45619$. В $\R^7$ задача уже решена
(теорема~\ref{thm:r7ex}): точное перечисление вершин ячейки с доказательством
полноты по $1$-скелету --- без предположения простоты многогранника --- дало
$1029$ статус \textbf{[С]}. Тот же инструмент применим и к $7203$/$28812$
дословно; препятствие там не в идее, а в размере, и мы его измерили. У решётки
индекса $7203$ в $\R^9$ ячейка имеет $752$ фасеты (потолок $2(2^9-1)=1022$) и
$1\,654\,170$ вершин против $254$ и $30\,368$ в $\R^7$; одно только
перечисление в плавающей точке заняло $10$ минут.

Попытка довести $\R^9$ до конца прошла три шага и встала на четвёртом; сообщаем
и результат, и место остановки, поскольку последнее оказалось не там, где
ожидалось.

\emph{Пройдено.} Поиск вороновски существенных векторов снимается априорной
границей: у релевантного вектора середина лежит на границе ячейки, поэтому
$\|v\|^2\le4R^2$, а строгую оценку $R$ сверху даёт лемма~\ref{lem:P1}
($4R^2\le4R_{\text{base}}^2+t^2=2929/400$). Класс чётности, не встретившийся в
этом окне, релевантного вектора не даёт: его кратчайший представитель длиннее, а
релевантный обязан быть кратчайшим в классе. Вместе с приведением базиса это
сжимает перебор с ${\sim}1{,}2\cdot10^{10}$ векторов (граница по
$\{0,1\}$-представителям при неприведённом базисе) до $4460$; на $\R^7$ тот же
приём даёт \emph{тот же} набор из $127$ пар, перебрав $486$ векторов вместо
$300\,368$. Сертификация всех $1\,218\,354$ вершин заняла $25$ минут --- этот
шаг узким местом не является.

\emph{Место остановки --- проверка полноты.} Она перебирает у каждой вершины все
$(N-1)$-подмножества активных фасет, что верно для простого многогранника, где
их ровно $N$. Здесь же есть костяк из ${\approx}4\,600$ вершин с $15$--$18$
активными фасетами, а $\binom{15}{8}=6435$: \emph{$0{,}28\,\%$ вершин порождают
$68\,\%$ подмножеств} --- $31{,}3$ миллиона из $46{,}1$. Существенно, что этот
костяк --- свойство самого ламинирования, а не выбора рациональной точки: у
симметричной рационализации (знаменатель Грама $40$) вершин на $26\,\%$ меньше,
$1\,218\,354$, но вырожденных ровно столько же, $4\,642$, и подмножеств даже
больше, $57{,}6$ миллиона. Значит, упирается не размерность, не число вершин и
не величина целых, а способ находить рёбра: для вырожденной вершины нужен
метод двойного описания, полиномиальный по числу самих экстремальных лучей, а
не перебор подмножеств. Альтернатива, обходящая вершины целиком: точное $R$ нам
и не нужно --- достаточно доказать $R^2<7/4$, то есть пустоту
$V_0\cap\{\|x\|^2\ge7/4\}$, а это задача ветвления с точными границами.

Попутно выяснились две вещи, меняющие план работ.

\emph{Кусочный путь не короче прямого.} Естественно было думать, что
рационализовать надо именно кусочный сертификат (\S\ref{ssec:cert}): он живёт в
размерности \emph{базы}, то есть на единицу меньше. Прямой счёт это опровергает:
у сертификата для $7203$ получается $256$ невырожденных областей, а вершин в них
суммарно $3\,475\,893$ --- вдвое больше, чем у самой девятимерной ячейки, причём
одна область даёт $159\,662$ вершины. Выигрыш размерности съедается тем, что
каждый кусок --- пересечение двух сдвинутых ячеек, а такое пересечение вершин
имеет больше, чем каждая из них. Поэтому целить надо в прямое перечисление
ячейки $\R^9$, а не в куски.

\emph{Рационализовать надо координаты сдвига, а не сам сдвиг.} Сдвиг слоя
хранится в декартовых координатах; округление их до знаменателя $10^4$ рушит
конструкцию --- слоевой минимум $D^2$ падает с $7$ до $4{,}35$. Если же перейти
к координатам базы и округлять там, конструкция выдерживает даже очень грубый
шаг: при знаменателе $1/20$ матрица Грама девятимерной решётки получается со
знаменателем $40$ и целыми записями до $360$ --- \emph{мельче}, чем у
сертификата в $\R^7$ (знаменатель $10^4$, записи до $45\,000$). Все пороги при
этом держатся: $R^2=1{,}689$ при нужном ${<}7/4$, минимум по горизонтальным
векторам ровно $7$, по слоевым $7{,}0365$, ширина $d=1{,}0179$. Это существенно:
стоимость точной арифметики определяется величиной целых, а не размерностью,
так что арифметическая часть в $\R^9$ дешевле, чем в $\R^7$, и всё упирается
ровно в число вершин. Данные и подстановка --- \path{campaigns/dim9_7203_exact.py}.

Запасы у оставшихся двух оценок различны: у $7203$ сертифицированный диаметр
отстоит от порога $\sqrt7$ на $1{,}6\,\%$, у $28812$ --- на $4{,}2\,\%$.
Полезный урок $\R^7$: рационализованная точка не обязана быть хуже численного
оптимума --- у неё связывающий слоевой запас оказался в тридцать с лишним раз
больше ($7{,}6\cdot10^{-5}$ против $2{,}3\cdot10^{-6}$).
\item \emph{Мозаичные рекорды Партса} \cite{Parts2023} ($k=8$: $1{,}4442$;
$k=15$: $2{,}3470$) ждут точных сертификатов --- машинерия
\S\ref{ssec:quant} применима к ним напрямую.
\end{itemize}
\end{enumerate}


